\section{Computation normalization and the raw comparison}
\label{sec:normalization}

This appendix states the inherited normalization argument explicitly, since
moving computation into syntax must count both its algorithm and its
metatheoretic justification~\cite{normalization,prior}. Reduction is the
contextual closure, including inside update bodies, of
\[
 I(0,a;u.h)\to a,\qquad I(Sn,a;u.h)\to h[I(n,a;u.h)/u].
\]
Capture-avoiding substitution preserves reduction. For a strongly normalizing
finite term $t$, let $\nu(t)$ be the largest length of a reduction sequence
from $t$, and let $|t|$ be its syntax size. The maximum exists because the
finitely branching reduction tree is well founded.

\begin{lemma}[Iterator closure]
If $n,a$ are strongly normalizing and $h[v/u]$ is strongly normalizing for
every strongly normalizing $v$, then $I(n,a;u.h)$ is strongly normalizing.
\end{lemma}
\begin{proof}
The body is strongly normalizing by taking $v=u$. Induct lexicographically
on $(\nu(n),|n|,\nu(a)+\nu(h))$. A contraction within $n$ lowers the first
coordinate; one in $a$ or $h$ lowers the third. Substitution stability
preserves the universal hypothesis on a reduced body. The zero root reduct
is $a$. At a successor root $n=Sp$, the counter $p$ has the same reduction
height and smaller size, so the induction hypothesis makes $I(p,a;u.h)$
strongly normalizing. The premise about $h$ makes its substitution into
the successor reduct strongly normalizing. Thus every immediate reduct is
strongly normalizing.
\end{proof}

\begin{theorem}[Normalization and substitution coherence]
Every finite raw term is strongly normalizing and has a unique normal form,
up to bound-variable renaming. With de Bruijn indices uniqueness is literal.
For every substitution $\sigma$,
\[
 \nf(t[\sigma])=
 \nf\bigl(\nf(t)[x\mapsto\nf(\sigma x)]\bigr).
\]
The same property holds for formulas.
\end{theorem}
\begin{proof}
Prove by structural induction that replacing a fixed term's free variables
by any strongly normalizing terms yields a strongly normalizing term. At an
iterator, the structural hypotheses for its counter, seed, and body are
exactly the preceding lemma's premises. Identity substitution gives strong
normalization. Local confluence follows because the two root patterns do
not overlap, disjoint contractions commute, and a contraction within a root
redex has finitely many residual copies after that root contraction. Contract
every retained copy; erased copies require no step. Substitution stability
gives a common reduct. Well-founded induction on reduction height now gives
confluence. Apply confluence to reducing $t$ before substitution and reducing
the supplied arguments before substitution to obtain the displayed identity.
Formula constructors themselves do not contract, so the formula statement
follows structurally.
\end{proof}

The delivered normalizer first traverses the children. At a successor
counter it normalizes the predecessor iterate and inserts that result into
the update by capture-avoiding substitution. Its contractions are precisely
the stated reductions, interspersed with finite structural traversals.
The theorem identifies its result with the unique normal form and justifies
termination in unbounded finite memory. Section~\ref{sec:bounded} separates
this total-recursive algorithm from the optional bounded certificate predicate.

\subsection{A computation inside an update is provably sound internally}

In the raw Choice-DI calculus of Appendix~\ref{sec:choice-introduction},
equation induction and ordinary equality reasoning are derived. If $u$ is
fresh for $\Gamma$ and $\Gamma\vdash h(u)=k(u)$, then
\[
 \Gamma\vdash I(n,a;u.h)=I(n,a;u.k).
\]
For this update-extensionality lemma, induct on the equation comparing the
two iterates at a fresh counter $z$. At zero both are $a$. In the successor
step, instantiate the supplied pointwise equality at the first iterate,
then use the induction equation inside $k$ to identify the two updates.
Only literal-equation induction is used. This lemma matters because a
free-parameter transport rule alone would not justify capture of that
parameter by an iterator binder.

Every contextual computation step is consequently an internal equality:
the root equations are axioms, usual congruence covers successor, counter and
seed positions, and the lemma covers update bodies. Thus every fixed term
satisfies $t=\nf(t)$ internally, and formula transport gives
$A\dashv\vdash\nf(A)$.

\subsection{The zero-witness case repair}
\label{sec:zero-repair}

Full normalization removes nonliteral closed aliases for zero. The following
inherited compound-envelope argument repairs a case-introduction proof whose
witness normalizes to zero~\cite{normalization,prior}. It is a derived
construction, not an extra primitive rule.

For an equation $E=(s=t)$ define
\[
 G_vE=(C(v,0,s)=C(v,0,t)),\qquad L(E)=q_vG_vE.
\]
The endpoints are literally $G_0E=(0=0)$ and $G_{Sw}E=E$.
The closure rules of Lemma~\ref{lem:row-cases} give $E\vdash L(E)$ and
$L(E)\vdash E$, using reflexivity and specialization at $1$ for the second.

Suppose a supplied derivation proves
$\Gamma,E(0,0)\vdash E(0,Sz)$. Put $A(x,y)=L(E(x,y))$.
Under $A(0,0)$, unwrap the envelope, insert that derivation by ordinary
assumption cut, and wrap its conclusion. Intro at witness zero now uses
cases, because $A$ is compound, and produces $D_{x,y}A(x,y)$.
Open it at fresh $c$. To introduce $D_{x,y}E(x,y)$ at witness $c$, assume
$E(c,0)$, wrap it as $A(c,0)$, apply the opened clause at $Sz$, and unwrap
$A(c,Sz)$. Since $c$ is not literal zero, this introduction again uses cases.
Closing the ordinary elimination proves the required $D_{x,y}E(x,y)$.
Thus unrestricted case introduction is available even at a zero witness.

\subsection{Exact raw/canonical comparison}

Normalize the terms of a raw Choice-DI derivation into an auxiliary canonical
calculus with separate case and equation-progress introductions. Computation
equations become reflexivity and normalized substitution coherence supplies
the remaining instances. The preceding repair supplies arbitrary case
introduction in the delivered calculus.

For arbitrary witnessed equation progress at a term $t$, first introduce
$D_{w,y}E(t,y)$ with unused $w$ at witness zero. This uses precisely the
progress hypothesis supplied. Open its conditional closure, use it under
$E(t,0)$ to prove the successor instance, and then use derived case
introduction to produce $D_{x,y}E(x,y)$ at $t$. Equation transport is
provided by the preceding new-rule theorem; replacing a fixed reflexivity
rule by instances of $\mathcal R$ changes no consequence.

Conversely, the raw calculus simulates the delivered rules using ordinary
paired elimination, its derived equality transport, and the appropriate case
or equation-progress introduction. Internal normalization equivalences
convert each canonical instance to its raw counterpart and back. The
canonical-to-raw-to-canonical composite is literal identity. The other
composite is the internally proved $A\dashv\vdash\nf(A)$. This proves the
raw/canonical comparison without assuming that a nonzero raw witness remains
nonzero after normalization, or that every rule tag survives arbitrary
numerical substitution.

\clearpage
\section*{Attribution and role of the following foundations}
\addcontentsline{toc}{section}{Attribution and role of the following foundations}
The next five appendices reproduce the logical and arithmetic reference
development from the pinned ProveIt sources~\cite{choice,erasure,forward,progress},
as reproduced in the preceding witness-discharge report~\cite{prior}.
They are licensed under the included MIT No Attribution license. They
establish the complete logical interface, predecessor, reflected equality,
quantifier-free induction, pairing, primitive-recursion compilation and both
interpretation composites. They introduce no additional primitive rule into
either delivered checker. Their historical labels ``old'' and ``new'' refer
to the inherited comparisons, not to Row in the present report.
